% Revize podle pozadavku ucitele, 2026-11-01.
% Nahrada dosavadni lekce o dusledkovych upravach: priprava z algebry,
% inspirace Realisticky.cz 020203-020204, ulohy ze zadanych fotografii.
\documentclass[11pt,a4paper]{article}
\input{../../sablony/pripravy/priprava.sty}
\nastavPripravu{1. ročník}{Rovnice a nerovnice}{73}{Algebraická příprava: složené závorky}
\begin{document}
\blok{Proč tuto hodinu zařazujeme}
V hodině 72 jsme zavedli kořen rovnice a ekvivalentní úpravy. Nyní procvičíme algebraické úpravy potřebné pro složitější rovnice: rozlišování kulatých, hranatých a složených závorek, násobení výrazů a kontrolu. Novou teorii rovnic nezavádíme; záměrně se neprobírají důsledkové úpravy.
\blok{Opakování postupu}
Před vnějšími závorkami nejdříve upravíme \emph{vnitřní} výraz. Po každém roznásobení kontrolujeme znaménka, zejména záporný násobitel. Zbývající rovnici dořešíme ekvivalentními kroky; výsledný kořen ověříme dosazením do původního zadání.
\begin{prikladReseniBox}[title={Příklad 1 – závorky ve třech úrovních}]
Vyřešte $2\{3[x-2(1-x)]+1\}=26$.
\tcblower
\begin{align*}
2\{3[x-2+2x]+1\}&=26\\
2\{3(3x-2)+1\}&=26\\
2(9x-5)&=26\\
18x-10&=26\\
18x&=36\\
x&=2.
\end{align*}
Zkouška v původním zadání: $2\{3[2-2(1-2)]+1\}=2(12+1)=26$.
\end{prikladReseniBox}
\begin{prikladReseniBox}[title={Příklad 2 – úloha podle fotografie (č. 1a)}]
Vyřešte $5\{5[5(5x-4)-4]-4\}=5$.
\tcblower
\begin{align*}
5\{5[25x-20-4]-4\}&=5\\
5\{5(25x-24)-4\}&=5\\
5(125x-124)&=5\\
625x-620&=5\\
625x&=625\\
x&=1.
\end{align*}
Přímá kontrola: $5\{5[5(5-4)-4]-4\}=5\{5(1)-4\}=5$.
\end{prikladReseniBox}
\poznamka{Na tabuli zapisuj celý člen při roznásobování, nezapisuj pouze hotovou hodnotu $x$. Příklad 2 je z učitelem dodané sbírky, nikoli z Realisticky.cz.}
\blok{Samostatné procvičování – stupňovaná obtížnost}
\begin{zadaniReseniBox}
\begin{enumerate}[label=\alph*)]
\item $4[2x-(x-3)]=20$.
\item $3\{2[3x-(x+1)]-4\}=12$.
\item $3[2(3x-6)-2(4x-5)+1]-3=6[3-8(x-3)]$ (úloha č. 1b z fotografie).
\end{enumerate}
\tcblower
\begin{enumerate}[label=\alph*)]
\item \(\begin{aligned}4(x+3)&=20\\x+3&=5\\x&=2.\end{aligned}\)
\item \(\begin{aligned}3\{2(2x-1)-4\}&=12\\3(4x-6)&=12\\12x-18&=12\\x&=\tfrac52.\end{aligned}\)
\item \(\begin{aligned}3[-2x-1]-3&=6[27-8x]\\-6x-6&=162-48x\\42x&=168\\x&=4.\end{aligned}\)
\end{enumerate}
\end{zadaniReseniBox}
\blok{Volitelná rezerva – skutečně náročná rovnice}
\begin{zadaniReseniBox}
Vyřešte $7\{x-2[5+x-8(x+3)-1]+2\}=3$.
\tcblower
\begin{align*}
7\{x-2[-20-7x]+2\}&=3\\
7(15x+42)&=3\\
105x+294&=3\\
105x&=-291\\
x&=-\frac{97}{35}.
\end{align*}
Zkouška lze provést dosazením do původní rovnice; zjednodušený levý výraz je $105x+294$.
\end{zadaniReseniBox}
\blok{Domácí úkol}
\begin{zadaniReseniBox}
\begin{enumerate}[label=\alph*)]
\item $5\{2[x-(3-x)]+1\}=35$.
\item $4[3-(2x-5)]=20$.
\item $x-(2x-0{,}2)^2=5{,}8-(2x+0{,}1)(2x-0{,}1)$ (úloha č. 1c ze sbírky).
\end{enumerate}
\tcblower
\begin{enumerate}[label=\alph*)]
\item \(\begin{aligned}5(4x-5)&=35\\4x-5&=7\\x&=3.\end{aligned}\)
\item \(\begin{aligned}4(8-2x)&=20\\32-8x&=20\\x&=\tfrac32.\end{aligned}\)
\item \(\begin{aligned}x-(4x^2-0{,}8x+0{,}04)&=5{,}8-(4x^2-0{,}01)\\1{,}8x-0{,}04&=5{,}81\\1{,}8x&=5{,}85\\x&=\tfrac{13}{4}.\end{aligned}\)
\end{enumerate}
\end{zadaniReseniBox}
\poznamka{Čas: připomenutí 5 min, první příklad 10 min, druhý 10 min, samostatně 15 min, závěr 5 min. Rezerva a domácí příklady se do 45 minut nepočítají. V 74. hodině následuje systematické řešení lineárních rovnic.}
\end{document}
